\section*{УВОД}
\label{sec:introduction}
\addcontentsline{toc}{section}{УВОД}

Разработката на графичен потребителски интерфейс за няколко операционни системи
едновременно поставя задача, чието обичайно решение се приема без проверка: всеки от
съществуващите многоплатформени инструменти твърди, че едно описание на интерфейса
стига за еднакво поведение върху различни платформи. Проверката на това твърдение е
трудна, защото за разлика от изчислителния код, чиято коректност се доказва с тест,
„правилният“ вид на потребителски интерфейс не е записан никъде във вид, който машина
може да сравни. Това е частен случай на \term{проблема за оракула} --- невъзможността
да се отличи желаното от нежеланото поведение при даден вход --- разгледан систематично
в~\cite{barr2015oracle}.

Настоящата работа изхожда от наблюдението, че проверката не е принципно невъзможна, а
само лишена от еталон. Ако съществува зряла реализация, приета за коректна, тя може да
служи не като образец за подражание, а като \term{измервателен инструмент} --- източник
на еталонни изображения, спрямо които независима реализация на същото описание се
сравнява количествено. Смяната на ролята е принципна: съпоставяната реализация
престава да бъде „копие“ и става второ, независимо решение на една и съща задача, а
разликите между двете стават измеримо свойство на абстракцията, а не въпрос на мнение.
Така изследването се отличава от инженерен отчет за пренасяне на код между езици: то
проверява хипотеза, която без независим втори изпълнител на едно и също описание би
останала недостъпна за проверка.

\textbf{Целта} на дипломната работа е да установи доколко наблюдаваното поведение на
графичен потребителски интерфейс може да бъде описано независимо от платформения
инструментариум, който го изобразява, и къде точно тази независимост се нарушава ---
чрез работеща реализация на такова описание върху няколко разнородни native
инструментариума, методика за количествена проверка на еквивалентността на
изобразяването и емпирична оценка на цената на избраната стратегия.

За постигането ѝ са формулирани следните \textbf{задачи}:

\begin{enumerate}[topsep=0pt,itemsep=0pt,leftmargin=*]
\item литературен преглед на подходите за многоплатформено описание на интерфейс и
      анализ на архитектурата, върху която стъпва изследването;
\item анализ на проектните алтернативи --- целеви език, стратегия за изграждане, модел
      на собствеността, набор от платформени слоеве --- и обосновка на избора;
\item проектиране на слоеста архитектура, в която описанието на интерфейса,
      изчисляването на разположението и системата от свойства са преносими, а
      платформено обусловената част е сведена до таблично зададено съответствие;
\item програмна реализация на архитектурата върху четири native инструментариума и
      един слой без екран;
\item методика и инструментариум за автоматизирано измерване на еквивалентността на
      изобразяването между две независими реализации на едно и също описание;
\item експериментална проверка върху пълния набор от интерфейси, систематизация на
      установените разминавания и оценка на цената на реализацията.
\end{enumerate}

\textbf{Обхватът} на работата е логическият слой, който описва интерфейса, изчислява
разположението и управлява свойствата, заедно с платформените слоеве, необходими за
неговото изобразяване върху избраните native инструментариума. Извън обхвата остават
съвместимостните слоеве към предходни поколения на еталонната реализация, вътрешните ѝ
помощни пространства от имена и допълненията, които не са част от ядрото ѝ.

\textbf{Методологична бележка.} Хипотезите за цената на реализацията --- размер на
артефакта, потребена памет, време за реакция --- са \term{регистрирани предварително},
преди да съществува измерване, заедно с изискванията за валидност, на които всяко
измерване трябва да отговаря, за да остане погрешна прогноза видима и отчетена като
такава, а не реконструирана след факта. Регистрацията е приложена към работата.

\textbf{Структура на изложението.} Раздел~\ref{sec:literature} разглежда предметната
област и архитектурния шаблон, върху който стъпва изследването.
Раздел~\ref{sec:alternatives} излага разгледаните проектни алтернативи и обосновава
избора между тях. Раздел~\ref{sec:design} представя проектирането на системата, а
Раздел~\ref{sec:implementation} --- нейната програмна реализация.
Раздел~\ref{sec:method} описва методиката за измерване на еквивалентността ---
самостоятелен принос на работата и предпоставка за валидността на всичко, което следва.
Раздел~\ref{sec:results} представя резултатите --- систематизацията на установените
нарушения на независимостта и измерената цена на реализацията. Заключението обобщава
приносите, ограниченията и насоките за по-нататъшна работа.
